% paper/sections/08d_defect_correction.tex
%
% §8.3  欠陥補正法（Defect Correction Method）による PPE 求解
%   — §8.2 で構築した CCD 離散演算子 L_H を高精度・低コストで解く枠組み
%   — Thomas 三重対角係数の導出：付録 \ref{app:sweep_thomas_coeff}
%   — 収束理論・スペクトル解析：付録 \ref{app:dc_convergence_theory}

\subsection{欠陥補正法による効率的 PPE 求解}
\label{sec:defect_correction_main}

§\ref{sec:ccd_poisson_matrix} で構築した CCD 離散演算子 $L_H$（$\Ord{h^6}$）は
高精度であるが，その\textbf{直接反復は困難}である．
$L_H$ は CCD ブロック Thomas 法（§\ref{sec:CCD}）に基づく密結合演算子であり，
(i)~帯幅が通常の中心差分（3点ステンシル）に比して広く，
(ii)~$x$-$y$ 方向の結合が強いため ADI 的な方向分離が自然には成立せず，
(iii)~大域疎行列を陽に構成して直接法（LU 分解）を適用すると
メモリ $\mathcal{O}(N^2)$ が必要となり大規模格子では非実用的である．

この困難を解消する枠組みが\textbf{欠陥補正法（Defect Correction Method）}である
\cite{Stetter1978}．
核心は「\textbf{高精度で評価し，低精度で解く}」という演算子分離にある：
右辺の残差（欠陥）には CCD 演算子 $L_H$（$\Ord{h^6}$）を用いて精度を確保し，
左辺の線形系には方向分離が容易な低次演算子 $L_L$（2次精度 FD）を用いて
$\mathcal{O}(N)$ コストで求解する．

\subsubsection{アルゴリズム}
\label{sec:dc_algorithm}

離散化された PPE $L_H\,p = b$ に対し，
初期推定 $p^{(0)} = 0$ から以下の3ステップを反復する：
\begin{equation}
  \boxed{%
  \begin{aligned}
    &\textbf{Step 1（欠陥の算出）：}\quad
      d^{(k)} = b - L_H\,p^{(k)}
    \\[4pt]
    &\textbf{Step 2（補正の求解）：}\quad
      L_L\,\delta p^{(k+1)} = d^{(k)}
    \\[4pt]
    &\textbf{Step 3（解の更新）：}\quad
      p^{(k+1)} = p^{(k)} + \omega\,\delta p^{(k+1)}
  \end{aligned}
  }
  \label{eq:dc_three_step}
\end{equation}

各記号の定義：
\begin{itemize}
  \item $L_H$：CCD 法で離散化した可変係数 Laplacian（$\Ord{h^6}$，§\ref{sec:ccd_poisson_matrix}）．
        行列不要で評価可能（各行・列に CCD ブロック Thomas 法を適用）．
  \item $L_L$：2次精度中心差分（FD）で離散化した近似 Laplacian．
        三重対角構造を持ち，$x$・$y$ 方向に\textbf{方向分離（directional sweep）}が可能．
  \item $d^{(k)}$：欠陥（defect）— 現在の近似解 $p^{(k)}$ を
        高精度演算子 $L_H$ で評価した残差．
  \item $\delta p^{(k+1)}$：補正量（correction）— 低次演算子 $L_L$ の逆作用で得る．
  \item $\omega$：緩和係数（relaxation parameter）—
        仮想時間刻み $\Delta\tau$ と数学的に等価（§\ref{sec:dc_terminology}）．
\end{itemize}

\subsubsection{方向分離スイープと Thomas 法}
\label{sec:dc_sweep}

Step~2 の補正方程式 $L_L\,\delta p = d$ において，
$L_L$ を $x$・$y$ 方向の1次元演算子に因子分解する：
\begin{equation}
  L_L = L_{L,x} + L_{L,y}
  \label{eq:ll_split}
\end{equation}
ここで $L_{L,x}$，$L_{L,y}$ はそれぞれ $x$・$y$ 方向の
変密度 FD 三重対角演算子（式~\eqref{eq:lfd_x_disc}）である．
この分離により，補正 Step~2 は2段の\textbf{方向分離スイープ（Directional Sweep）}に帰着する：
\begin{align}
  \text{$x$ スウィープ：}\quad
    \bigl[\tfrac{1}{\Delta\tau}\,\mathbf{I} - L_{L,x}\bigr]\,q &= d^{(k)}
    \label{eq:dc_sweep_x} \\
  \text{$y$ スウィープ：}\quad
    \bigl[\tfrac{1}{\Delta\tau}\,\mathbf{I} - L_{L,y}\bigr]\,\delta p^{(k+1)} &= q
    \label{eq:dc_sweep_y}
\end{align}
各スウィープは独立な三重対角系のバッチであり，
Thomas 法（付録~\ref{app:sweep_thomas_coeff}）により $\mathcal{O}(N)$ で求解される．
$x$ スウィープでは $N_y$ 本の $N_x$ 節点三重対角系が，
$y$ スウィープでは $N_x$ 本の $N_y$ 節点三重対角系が，
それぞれベクトル化された並列 Thomas 法で一括処理される．
1反復あたりの総計算量は $\mathcal{O}(N^2)$（$N = N_x = N_y$）であり，
Krylov 法（GMRES 等）の1ステップと同オーダーだが Thomas 法の定数係数が小さく実測コストが低い．

\subsubsection{収束先の精度保証}
\label{sec:dc_accuracy}

欠陥補正法の本質的な性質は，
\textbf{左辺に低次演算子 $L_L$ を用いても，収束先の解の空間精度は損なわれない}点にある．
収束条件 $\|d^{(k)}\| < \varepsilon_\mathrm{tol}$ を満たした時点で，
\begin{equation}
  L_H\,p^{(k)} \approx b
  \quad\text{（精度 $\varepsilon_\mathrm{tol}$）}
\end{equation}
が成立し，最終圧力解は CCD 演算子 $L_H$ の $\Ord{h^6}$ 精度を完全に保持する．
$L_L$ は収束を駆動する「加速装置」の役割のみを担い，
収束先の精度は右辺の欠陥評価（$L_H$ による残差）が決定する．

この精度保証は反復回数 $k$ と密接に関連する：
$k=1$ では $L_L$ の $\Ord{h^2}$ が支配的だが，
$k$ を増やすごとに $L_H$ の高次精度が発現し，
$k \geq 3$ で $\Ord{h^6}$ に漸近する
（定量的検証は §\ref{sec:dc_iteration_accuracy} に示す）．
反復行列のスペクトル解析（付録~\ref{app:dc_convergence_theory}）により，
$\omega$ の最適値と収束速度が理論的に導かれる．

\subsubsection{PPE 求解法の比較}
\label{sec:dc_comparison}

表~\ref{tab:dc_ppe_methods} に本稿で使用可能な PPE 求解法を比較する．
欠陥補正法を標準ソルバとして採用し，
LGMRES/直接 LU はデバッグ・検証用途に限定する．

\begin{table}[htbp]
\centering
\small
\caption{PPE 求解法の比較}
\label{tab:dc_ppe_methods}
\renewcommand{\arraystretch}{1.3}
\begin{tabular}{lcccc}
\toprule
手法 & 空間精度 & 行列構築 & 1反復コスト & 用途 \\
\midrule
BiCGSTAB（FVM） & $\Ord{h^2}$ & 必要 & $\mathcal{O}(N^2)$ & テスト専用 \\
LGMRES + CCD Kronecker & $\Ord{h^6}$ & 必要 & $\mathcal{O}(N^2)$ & 検証用 \\
直接 LU + CCD Kronecker & $\Ord{h^6}$ & 必要 & $\mathcal{O}(N^3)$ & 検証用 \\
\textbf{欠陥補正法（本稿標準）} & $\boldsymbol{\Ord{h^6}}$ & \textbf{不要} & $\boldsymbol{\mathcal{O}(N^2)}$ & \textbf{本番} \\
\bottomrule
\end{tabular}

\smallskip
\noindent 注：欠陥補正法は右辺残差に $L_H$（CCD, $\Ord{h^6}$），
左辺に $L_L$（FD, $\Ord{h^2}$）を用いる（§\ref{sec:dc_algorithm}）．
行列不要（matrix-free）で $\mathcal{O}(N^2)$/反復．
Thomas 係数の導出は付録~\ref{app:sweep_thomas_coeff}，
収束理論は付録~\ref{app:dc_convergence_theory} を参照．
\end{table}

\subsubsection{用語の統一と対応関係}
\label{sec:dc_terminology}

本稿における欠陥補正法関連の用語を以下に整理する．

\begin{table}[htbp]
\centering
\small
\caption{欠陥補正法の用語定義と対応関係}
\label{tab:dc_terminology}
\renewcommand{\arraystretch}{1.3}
\begin{tabular}{lll}
\toprule
用語 & 定義 & 備考 \\
\midrule
欠陥（Defect） & $d^{(k)} = b - L_H\,p^{(k)}$ &
  高精度演算子 $L_H$ で評価した残差 \\
補正（Correction） & $\delta p^{(k+1)}$：$L_L\,\delta p = d$ の解 &
  低次演算子 $L_L$ の逆作用 \\
方向分離スイープ & $x$/$y$ 方向の逐次 Thomas 法 &
  $L_L$ の因子分解実行手順 \\
  （Directional Sweep） & （式~\eqref{eq:dc_sweep_x}--\eqref{eq:dc_sweep_y}） & \\
仮想時間刻み & $\Delta\tau$：$1/\Delta\tau$ は $L_L$ の対角シフト &
  欠陥補正の緩和係数 $\omega$ と \\
  （Virtual Time Step） & （式~\eqref{eq:dc_sweep_x}） &
  数学的に等価$^{\dagger}$ \\
\bottomrule
\end{tabular}

\smallskip
\noindent $^{\dagger}$
式~\eqref{eq:dc_three_step} の $\omega$ と仮想時間法の $\Delta\tau$ は
$[\frac{1}{\Delta\tau}\mathbf{I} - L_L]\,\delta p = d$ を $\delta p = \Delta\tau\,[\mathbf{I} - \Delta\tau L_L]^{-1} d$ と整理することで
$\omega = \Delta\tau$ の関係にあることが確認できる
（厳密な導出は付録~\ref{app:dc_convergence_theory}）．
\end{table}
